%=====================================================================
\chapter{Introduction}
\label{ch:intro}
%=====================================================================

\section{Stock Market Decision Making as an Information Problem}

The price of a traded security is the outcome of a continuous negotiation between
participants who do not observe the same information at the same time. Part of that
information is structured and numerically clean: the open, high, low, close and volume
records of each session, and the technical indicators derived from them. A larger part is
unstructured and arrives irregularly, in news articles, corporate disclosures, regulatory
announcements, analyst commentary and social media discussion. A third part changes slowly
and shapes the environment rather than the individual security, in the form of policy rates,
inflation, industrial output, currency movements and commodity prices. Any forecasting
system that reads only one of these streams is, by construction, reasoning from an
incomplete description of the market it is trying to
anticipate~\cite{snasel2024,long2024,nti2021}.

This is not a new observation. What is new is that acting on it has become technically
feasible. Two developments are responsible. The first is the maturation of multi-source
data fusion, which supplies a principled account of how heterogeneous streams may be
combined so that their complementary content is preserved instead of being averaged
away~\cite{snasel2024,zhang2022fusion,li2023datafusion}. The second is the arrival of
transformer-based large language models, which read financial text with a level of
contextual sensitivity that lexicon and bag-of-words methods never achieved, and which can
therefore convert narrative information into a numerically usable
signal~\cite{vaswani2017,araci2019,li2024llmsurvey}. Together they make it reasonable to ask
a question that was previously out of reach: whether a single system can absorb
quantitative, textual and macroeconomic evidence and convert it into a decision that an
investor could actually execute.

The framing matters. A very large share of the published literature in this area treats the
problem as one of \emph{prediction}: given a feature vector, forecast tomorrow's return or
its sign, and report an accuracy figure. That framing is legitimate but incomplete, because
an investor does not consume a prediction. An investor consumes a position --- a vector of
holdings, sized against a risk budget, rebalanced at a cost, and held through drawdowns.
The translation from the first object to the second is not a formality. It involves a
covariance structure that must be estimated, position limits that bind, turnover that must
be paid for, and a decision about how much total risk to carry that is logically separate
from the decision about what to hold. Each of those steps can destroy the value of a
forecast as easily as preserve it~\cite{demiguel2009,gu2020,markowitz1952}.

\section{Three Obstacles Between a Fused Signal and a Defensible Decision}

Three difficulties sit between the raw heterogeneous information and a defensible decision,
and this thesis is directed at all three simultaneously rather than at any one of them in
isolation.

\subsection{The reliability of a source is not constant}

Technical indicators are informative in quiet trending conditions and degrade when the
market is disturbed; news and volatility signals behave in the opposite manner, carrying
little information in placid stretches and a great deal during a shock. A fusion rule with
fixed weights therefore misweights its inputs precisely when accurate weighting matters
most~\cite{farimani2024,tai2022}. This is a structural problem rather than a tuning problem.
Attention mechanisms, which are the standard remedy, learn weights over an entire training
set; the resulting weight expresses the \emph{average} relevance of a source across the
training period and cannot express the proposition that a particular source has become
unreliable in the last fortnight. Between retraining cycles, a stream that has degraded
continues to influence the forecast at its historical weight.

What is missing is a measurement. If the reliability of a source can be estimated
continuously from observable evidence --- specifically, from the variance of that source's
own recent prediction errors --- then the fusion weight becomes a function of that estimate
and the architecture self-corrects without human intervention and without waiting for the
next training run. Chapter~\ref{ch:dynfusion} develops this idea formally and tests it.

\subsection{A prediction is not an explanation}

A forecast produced by a deep architecture is not, by itself, an account of why the forecast
was made. In a domain where capital is committed on the basis of the output, and where a
regulator or an investment committee may reasonably demand a justification, opacity is a
practical obstacle and not merely a philosophical one~\cite{koa2024,lundberg2017}.

The standard response is post-hoc attribution: methods such as Shapley additive
explanations~\cite{lundberg2017} or local surrogate models~\cite{ribeiro2016} report which
inputs the fitted model relied upon. This is genuinely useful, and it is not what this thesis
pursues as its primary transparency instrument. The reason is that post-hoc attribution
explains a \emph{model}; it does not establish that the relationships the model has learned
are anything more than correlations that happened to hold in the training window. A feature
that co-moves with the target for spurious reasons will be attributed importance exactly as
a genuine driver would be. Chapter~\ref{ch:causal} therefore pursues transparency at the level
of \emph{structure}: the feature set entering the model is constrained, before any learning
takes place, to variables with demonstrable directed influence, recovered by conditional
independence testing over lagged series~\cite{spirtes2000,granger1969,xu2026causal}.

\subsection{Statistical accuracy and economic usefulness are different properties}

The third obstacle is the one most often elided. A model that is accurate is not thereby
profitable, and a model that is profitable is not necessarily accurate. The two properties
come apart for reasons that are well understood in principle and routinely ignored in
practice: predictions are not equally costly to act on, a small edge can be consumed entirely
by transaction costs, an unconditional accuracy figure says nothing about performance in the
states where losses are concentrated, and --- most importantly --- a portfolio's risk-adjusted
return depends on the timing of risk as much as on the selection of assets.

This thesis treats the decision layer as a first-class object of study rather than as a
postscript. Chapter~\ref{ch:allocation} develops a decision-theoretic allocation framework,
evaluates it over eighteen years with transaction costs, and --- unusually --- disassembles it
to determine which of its layers actually produce the measured outcome. The answer, reported
in full including the parts that do not favour the original design, is the central finding of
the thesis.

\section{Research Context and the Position of the Four Objectives}

Figure~\ref{fig:context} places the four research objectives on the path from heterogeneous
raw information to an executable decision. Objective~1 governs the integration of sources.
Objective~2 governs the extraction of signal from language and its combination with
conventional indicators. Objective~3 governs the transparency of the reasoning and the
behaviour of sentiment as a noisy, time-lagged input. Objective~4 governs the conversion of a
signal into a position and the validation of that conversion across markets.

\begin{figure}[!htb]
\centering
\scalebox{0.86}{\begin{tikzpicture}[node distance=3.2mm]
\node[blk,minimum width=23mm] (m) {Market data\\OHLCV, indicators};
\node[blk,right=of m,minimum width=23mm] (n) {News and social\\media text};
\node[blk,right=of n,minimum width=23mm] (mac) {Macroeconomic\\series};
\node[blk,right=of mac,minimum width=23mm] (alt) {Alternative data\\volatility index};

\node[blkg,below=9mm of n,minimum width=52mm,xshift=9mm] (fuse)
  {\textbf{Objective 1} --- multi-source integration\\
   time-varying source weighting};
\node[blka,below=7mm of fuse,minimum width=52mm] (llm)
  {\textbf{Objective 2} --- LLM sentiment extraction\\
   hybrid fusion with financial indicators};
\node[blka,below=7mm of llm,minimum width=52mm] (xai)
  {\textbf{Objective 3} --- transparent reasoning\\
   sentiment noise and temporal dynamics};
\node[blkg,below=7mm of xai,minimum width=52mm] (alloc)
  {\textbf{Objective 4} --- validated decision making\\
   allocation, risk budget, cross-market test};
\node[blkd,below=7mm of alloc,minimum width=52mm] (dec)
  {Executable position with a stated risk profile};

\draw[ar] (m.south) -- ++(0,-3mm) -| (fuse.north west);
\draw[ar] (n.south) -- ++(0,-3mm) -| ([xshift=-14mm]fuse.north);
\draw[ar] (mac.south) -- ++(0,-3mm) -| ([xshift=14mm]fuse.north);
\draw[ar] (alt.south) -- ++(0,-3mm) -| (fuse.north east);
\draw[ar] (fuse) -- (llm);
\draw[ar] (llm) -- (xai);
\draw[ar] (xai) -- (alloc);
\draw[ar] (alloc) -- (dec);
\node[lbl,right=5mm of fuse,align=left] {what to fuse};
\node[lbl,right=5mm of llm,align=left] {how to read text};
\node[lbl,right=5mm of xai,align=left] {why to trust it};
\node[lbl,right=5mm of alloc,align=left] {how much to hold};
\end{tikzpicture}}
\caption[Research context and the position of the four objectives]{Research context: the
path from heterogeneous market information to an executable decision, and the position of
the four research objectives on that path.}
\label{fig:context}
\end{figure}

The objectives are sequential in the sense that each consumes the output of the one before
it, but they are not merely stages of a pipeline. Each raises a distinct research question,
each is answered by distinct experimental evidence on a study universe chosen for that
question, and --- importantly for how this thesis should be read --- the answer obtained at
one stage redirected the design of the next. The framework reported in
Chapter~\ref{ch:allocation} is not the framework that was envisaged at the outset. It is the
framework that survived the evidence.

\section{The Indian Equity Market as the Primary Empirical Setting}

The Indian equity market is the primary empirical setting of this research. The choice is
deliberate rather than convenient, and it is worth stating the reasons explicitly because
they also bound the generality of several results.

The National Stock Exchange combines four features that make it an unusually demanding test
of a multi-source system. First, retail participation is high and has risen sharply, so
sentiment expressed on social platforms is not merely commentary on the market but a
partial cause of short-term price movement~\cite{khan2025asymmetric,mu2023investor}. Second,
the market is strongly sensitive to domestic policy events --- the annual Union Budget,
Reserve Bank of India monetary policy committee announcements, and regulatory decisions by
the Securities and Exchange Board of India --- which produce discrete structural breaks
rather than smooth transitions. Third, the index carries pronounced sectoral concentration
in financial services, so a shock in that sector propagates to the aggregate in a way it
would not in a more evenly weighted index. Fourth, the market is simultaneously exposed to
global factors, including crude oil prices, the rupee--dollar exchange rate and United States
monetary policy, so domestic and international information can and frequently do
disagree~\cite{bacco2024,dong2024}.

That disagreement is the point. A fusion architecture is only interesting when the sources
it fuses carry genuinely different information; in a market where every stream says the same
thing, the fusion rule is irrelevant and a single source would do. The Indian market supplies
sources that disagree, and the disagreement is informative rather than noise.

Two qualifications are recorded here and revisited where they bind. The first is that
instrument-level, time-stamped news corpora of the depth available for United States
large-capitalisation equities do not exist for broad Indian index instruments, which
constrains what can be measured at the allocation layer and is the direct reason the
sentiment construct used in Chapter~\ref{ch:allocation} is market-derived rather than
textual. The second is that the allocation study of Chapter~\ref{ch:allocation} requires
eighteen years of daily data on a universe with genuine dispersion in asset volatility,
including two severe crises; the instruments that satisfy that requirement with clean,
survivorship-free history are United States exchange-traded funds, and the study is
therefore conducted there and treated as a cross-market test of the mechanism rather than as
a claim about Indian allocation. Both decisions are stated rather than concealed, because
they determine precisely which claims in this thesis are about the Indian market and which
are about the mechanism.

\section{Research Questions}

The thesis addresses five research questions, each of which is answered by evidence reported
in a specific chapter.

\begin{enumerate}[label=\textbf{RQ\arabic*.},leftmargin=2.8em]
\item Does fusing financial, textual and macroeconomic streams under a learned, time-varying
weighting improve forecast quality over the best single-source and fixed-weight alternatives,
and can the improvement be attributed to individual streams rather than merely asserted?
(Chapter~\ref{ch:fusion}.)

\item Can the weight assigned to a source be made a function of that source's measured,
instantaneous predictive reliability rather than its historical average relevance, and does
doing so improve risk-adjusted outcomes? (Chapter~\ref{ch:dynfusion}.)

\item How much of the improvement that large language models deliver in sentiment
\emph{measurement} survives into forecast quality once the sentiment score is temporally
aligned, fused with technical indicators, and evaluated out of sample at the horizon and
granularity at which decisions are actually taken? (Chapter~\ref{ch:llm}.)

\item Can transparency be grounded in directed causal structure rather than in post-hoc
attribution, and can the noise carried by the sentiment channel be measured rather than
assumed? (Chapter~\ref{ch:causal}.)

\item When the fused signal is converted into an executable portfolio under realistic
constraints, which layers of the resulting framework actually produce the measured outcome,
and under what conditions does the framework transfer to a different market?
(Chapter~\ref{ch:allocation}.)
\end{enumerate}

\section{Contributions of the Thesis}

The contributions claimed are stated here in summary and developed with their supporting
evidence in the chapters indicated. They are stated as bounded claims; where the evidence
supports only a qualified statement, the qualification is part of the contribution.

\begin{enumerate}[label=\textbf{C\arabic*.},leftmargin=2.8em]
\item \textbf{An attributable four-stream fusion architecture for an emerging market.} A
Multi-Source Fusion Network integrating technical indicators, news sentiment, social media
sentiment and Indian macroeconomic series through modality-specific encoders and a
cross-modal attention layer, evaluated on ten years of National Stock Exchange data across
ten sectoral indices under a strictly chronological protocol with expanding-window
walk-forward confirmation. The contribution is not that fusion improves accuracy, which was
already known, but that the improvement is \emph{attributable}: source ablation isolates the
marginal value of each stream and demonstrates superadditivity between the two textual
sources, and the learned attention distribution demonstrates that source importance is
state-dependent. (Chapter~\ref{ch:fusion}.)

\item \textbf{A fusion rule conditioned on measured source reliability.} A Bayesian dynamic
uncertainty-weighted formulation in which each source is modelled by its own expert network,
its predictive uncertainty is estimated continuously from the variance of its recent errors,
and the fusion weights are obtained by a softmax over the negative uncertainties. The
contribution is isolated by construction: the same expert networks fused with fixed weights
perform no better than the established static baselines, which locates the measured gain in
the weighting mechanism rather than in the networks. (Chapter~\ref{ch:dynfusion}.)

\item \textbf{A measured account of what language-model sentiment contributes.} A hybrid
pipeline coupling a transformer sentiment encoder with fifteen technical indicators,
evaluated against technical-only baselines under a like protocol. Three findings constitute
the contribution: sentiment is the highest-gain feature in the fused space; it is orthogonal
to price, correlating 0.04 with price level and 0.74 with momentum; and the incremental gain
in raw directional accuracy is small, with all models close to the random classifier at a
single-name daily horizon. The third finding is reported as prominently as the first two, and
it is what redirects the thesis toward the decision layer. (Chapter~\ref{ch:llm}.)

\item \textbf{Structural rather than post-hoc transparency, with sentiment noise quantified.}
A Dynamic-Causal Hybrid Framework that constrains the feature set by directed causal structure
recovered through conditional independence testing, encodes intraday, daily and weekly
horizons through a multi-scale memory encoder, and allocates through a credibilistic
optimiser. It recovers an interpretable causal graph of density 0.850 with an identifiable
causal hub, and demonstrates that strong directed influence can coexist with weak linear
correlation --- the mechanistic justification for causal rather than correlational feature
selection. Separately, the sentiment construct is validated against four independently
published series and shown to separate identified market states at $p = 1.8 \times 10^{-8}$,
while its predictive content is measured independently and found not to be significant.
Keeping construct validity and predictive usefulness apart is itself a methodological
contribution. (Chapter~\ref{ch:causal}.)

\item \textbf{A decision framework ablated at matched risk, with its applicability bounded by
evidence.} A regime-conditional allocation framework that separates relative composition from
total risk, evaluated over eighteen years with transaction costs, a fully disclosed
specification sweep, four pre-selected stress episodes, a 665-session holdout postdating every
design decision, and a secondary market universe. The ablation is conducted \emph{at matched
average exposure}, which is what makes it possible to distinguish the timing of risk from its
level; it finds that regime-conditional risk timing improves both the Sharpe ratio and the
drawdown with no significant change in mean return, and that the signal fusion layer
contributes nothing detectable. (Chapter~\ref{ch:allocation}.)

\item \textbf{Two generalisable engineering failure modes.} Both were found only by
disassembling the framework and both are reported because they are silent: an absolute
volatility target calibrated for a concentrated equity portfolio binds in one rebalance out of
224 on a diversified universe, leaving the regime layer present in the specification but
inert; and a signal-to-allocation conversion in which mismatched units between the linear and
quadratic terms of a mean-variance objective, rather than the information in the signal,
determine the optimiser's solution. Neither produces an error message.
(Chapter~\ref{ch:allocation}.)
\end{enumerate}

\section{Scope and Delimitations}

Six delimitations bound the work and are stated at the outset so that no claim in later
chapters is read more widely than the evidence supports.

\textbf{Horizon.} All studies operate at a daily decision horizon. Intraday microstructure,
order-book dynamics and execution algorithms are outside scope. The multi-granularity encoder
of Chapter~\ref{ch:causal} uses intraday-derived features but produces daily decisions.

\textbf{Direction of exposure.} Every portfolio constructed in this thesis is long-only and
unlevered. Risk budgeting is one-sided: the framework de-risks toward cash in disturbed states
but never levers up. This is a deliberate restriction, because it guarantees that any measured
performance difference is attributable to risk reduction rather than to leverage.

\textbf{Language coverage.} Sentiment extraction is English-only. A substantial share of Indian
retail participants consume financial news in regional languages, and the sentiment channel
measured here therefore observes only part of the information actually reaching the market.
This is identified as future work in Chapter~\ref{ch:conclusion} rather than dismissed.

\textbf{Transfer.} The cross-market evidence in Chapter~\ref{ch:allocation} comes from applying
the framework to a second universe \emph{without modification}. This is direct transfer, not
transfer learning: no parameter was fine-tuned on the target market, no representation was
pre-trained and re-used, and no domain-adaptation objective was optimised. The result is
therefore a lower bound on cross-market performance, and it is labelled as such wherever it
appears.

\textbf{Explainability methods.} Transparency in this thesis rests on directed causal structure
and on the interpretability of the fusion weights. Attribution methods such as Shapley additive
explanations and local surrogate models are reviewed and cited but are not run as experiments;
where the thesis speaks of explainability it means structural transparency, and the distinction
is maintained throughout.

\textbf{Costs and frictions.} Transaction costs are modelled as a one-way charge of ten basis
points in the allocation study and are absent from the prediction studies of
Chapters~\ref{ch:fusion} to~\ref{ch:llm}, where the evaluation is of forecast quality rather
than of realised profit. Slippage, market impact and financing costs are not modelled. The
consequence --- that directional accuracy reported in Chapter~\ref{ch:fusion} does not
translate directly into trading profit --- is stated in that chapter rather than left implicit.

\section{Organisation of the Thesis}

The thesis is organised into nine chapters. Chapter~\ref{ch:background} presents the technical
foundations required to read the later chapters without recourse to external material:
technical indicator construction, the three generations of sentiment extraction, fusion
strategies, regime modelling, portfolio construction, and the evaluation metrics and
statistical procedures used throughout. Chapter~\ref{ch:litsurvey} reviews the research
literature along the four axes of the objectives and derives the five research gaps that
define the scope of the work. Chapter~\ref{ch:fusion} develops and evaluates the Multi-Source
Fusion Network. Chapter~\ref{ch:dynfusion} develops the Bayesian dynamic uncertainty-weighted
formulation of the fusion rule. Chapter~\ref{ch:llm} develops the hybrid language-model and
indicator framework. Chapter~\ref{ch:causal} develops the explainable and causal framework and
the sentiment diagnostics. Chapter~\ref{ch:allocation} develops the decision-theoretic
allocation framework and reports the ablation, the holdout and the cross-market test.
Chapter~\ref{ch:conclusion} consolidates the findings, states the limits of the evidence and
sets out the future scope. A statement of the publications underlying the thesis and the
consolidated reference list follow.
